\cvsection{OUTSIDE COURSEWORK}

	\cvevent{Quantum Computation Internship}{\href{https://www.iiserkol.ac.in/web/en/people/faculty/dps/pprasanta/}{Prof. Prasanta K. Panigrahi}}{June 2022 -- July 2022}{\href{https://www.iiserkol.ac.in/web/en/}{IISER Kolkata}}

	\hspace{5mm} Here I learned the basics of quantum computation and various aspects of it. I read a lot of papers suggested by my instructor and also did some activities on my own. Finally, I read a paper on a Quantum Robot and felt I could solve the problem addressed in the paper better. So I learned more about the topic and submitted my report on the same.

	\hspace{5mm} Internship Report: \href{https://github.com/PeithonKing/Quantum_Robot_LaTEX/blob/main/quantum_robot_internship_report.pdf}{\textbf{Quantum Robot}}

	\divider


	\cvevent{PyaR Seminar 2021}{\href{https://www.astro.ucsc.edu/faculty/index.php?uid=pguhatha}{Prof. Raja GuhaThakurata}}{July $29^{th}$ to $31^{st} 2021$}{Online}

	\hspace{5mm} Here we learned the python programming language and how it can be used with Jupyter Notebook. We also learned the basics of astronomical data analysis using libraries like Numpy, Pandas, matplotlib etc. We were also briefed about some clustering algorithms used daily in this field.

	\hspace{5mm} GitHub Repository (materials): \href{https://github.com/PeithonKing/PyaR-2021}{\textbf{PeithonKing/PyaR-2021}}

	\divider


	\cvevent{Quantum Computation Course}{\href{http://www.iisertirupati.ac.in/}{IISER Tirupati} & \href{https://qkrishi.com/}{Qkrishi}}{2022 Summer Break}{Online}

	\hspace{5mm} I did a course on the basics of Quantum Computation and Quantum Information jointly organised by \href{http://www.iisertirupati.ac.in/}{IISER Tirupati} and Qkrishi. We learned the basic theory and had a hands-on experience with the IBM Quantum Experience. We also submitted a term project of Attacking Quantum Key Distribution Protocols. We demonstrated QKD algorithms like BB84 and E91 protocols in multiple devices on the same network using python libraries like Flask and Qiskit.

	\hspace{5mm} GitHub Repository (materials): \href{https://github.com/PeithonKing/Attacking_QKD_Protocols}{\textbf{PeithonKing/Attacking\_QKD\_Protocols}}

	\divider
